\section{Introduction} \label{sec:introduction}
Catastrophic forgetting (CF) and knowledge transfer (KT) are two key challenges of continual learning (CL), which learns a sequence of tasks incrementally. CF refers to the phenomenon where a model loses some of its performance on previous tasks once it learns a new task. KT means that tasks may help each other to learn by sharing knowledge.
This work further investigates these problems in the popular CL paradigm, \textit{task-incremental learning} (TIL). In TIL, each task consists of several classes of objects to be learned. Once a task is learned, its data is discarded and will not be available for later use. 
During testing, the task id is provided for each test sample so that the corresponding classification head of the task can be used for prediction.  


Several effective approaches have been proposed for TIL that can achieve learning with little or no CF. \textit{Parameter isolation} is perhaps the most successful one in which the system learns to mask a sub-network for each task in a shared network.
HAT~\citep{hat} and SupSup~\citep{supsup} are two representative systems.
HAT set binary/hard masks on neurons (not parameters) that are important for each task. In learning a new task, those masks block the gradient flow through the masked neurons in the backward pass. Only those \textit{free} (unmasked) \textit{neurons} and their parameters are trainable. Thus, as more tasks are learned, the number of free neurons left becomes fewer, making later tasks harder to learn, which results in gradual performance deterioration (see \cref{sec:capacity}). Further, if a neuron is masked, all the parameters feeding to it are also masked, which consumes \textbf{a great deal of network capacity} (hereafter referred to as the ``capacity problem''). 
As the sub-networks for old tasks cannot be updated, it has \textbf{limited knowledge transfer}.  
CAT~\citep{cat} tries to improve KT of HAT by detecting task similarities. If the new task is found similar to some previous tasks, these tasks' masks are removed so that the new task training can update the parameters of these tasks for backward pass. However, this is risky because if a dissimilar task is detected as similar, serious CF occurs, and if similar tasks are detected as dissimilar, its knowledge transfer will be limited.
SupSup uses a backbone network (randomly initialized) and finds a sub-network for each task. The sub-network is represented by a mask, which is a set of binary gates indicating which parameters in the network are used. The mask for each task is saved. Since the network is not changed, SupSup has no CF or capacity problem, but since each mask is independent of other masks, SupSup by design has no KT. 


To tackle these problems, we propose a very different approach, named ``\underline{\textit{S}}oft-masking of \underline{\textit{P}}arameter-level \underline{\textit{G}}radient flow'' (SPG). It is surprisingly effective and contributes in following ways:

(1).~Instead of learning hard/binary masks on neurons for each task and blocking these neurons in training a new task and in testing like HAT, SPG computes an importance score for each network parameter (not neuron) to old tasks using gradients. The reason that gradients can be used as importance is because gradients directly tell how a change to a specific parameter will affect the output classification and may cause CF. SPG uses the importance score of each parameter as a \textit{soft-mask} to constrain the gradient flow in the backward pass to ensure those important parameters to old tasks have minimum changes in learning a new task to prevent CF of previous knowledge. 
To our knowledge, soft-masking of parameters has not been done before. 

(2).~SPG has some resemblance to the popular regularization-based approach, e.g., EWC~\citep{ewc}, in that both use importance of parameters to constrain changes to important parameters of old tasks. But there is a major difference. SPG directly controls each parameter (fine-grained), but EWC controls all parameters together using a regularization term in the loss to penalize the sum of changes to all parameters in the network (rather coarse-grained). \cref{sec:results} shows that our \textbf{soft-masking is markedly better than regularization}. We believe this is an important result.


(3).~In the forward pass, no masks are applied, which encourages knowledge transfer among tasks. This is better than CAT as SPG does not need extra mechanism for task similarity comparison. Knowledge sharing and transfer in SPG are automatic. SupSup cannot do knowledge transfer. 


(4).~As SPG soft-masks parameters, it does not monopolize any parameters or sub-network like HAT for each task and SPG's forward pass does not use any masks. This reduces the \textit{capacity problem}. 


Experiments with the standard CL setup have been conducted with (1) similar tasks to demonstrate SPG's better knowledge transfer 
and (2) dissimilar tasks to show SPG's ability to overcome CF, and (3) deal with the capacity issue. None of the baselines is able to achieve all.
The code is available at \url{https://github.com/UIC-Liu-Lab/spg}.


